\documentclass[11pt,a4paper]{article}
% Shared preamble for RuPhO English problem statements (match T1 2026 style).
% Use: \input{latex/rupho_en_preamble.tex} after \documentclass.
\usepackage[margin=1in]{geometry}
\usepackage{amsmath,amssymb}
\usepackage{graphicx}
\usepackage{float}
\usepackage{enumitem}
\usepackage{microtype}
\usepackage{caption}
\usepackage{tabularx}
\usepackage{hyperref}

\captionsetup{font=small,labelfont=bf,skip=6pt}

\setlist[itemize]{leftmargin=1.2cm}
\setlist[enumerate]{leftmargin=1.2cm}

% One outer box per subproblem: [id | statement | points] (no inner rules)
\newcommand{\problembox}[3]{%
  \par\addvspace{0.42em}%
  \noindent
  \setlength{\fboxsep}{7.5pt}%
  \setlength{\fboxrule}{0.95pt}%
  \fbox{%
    \begin{minipage}{\dimexpr\linewidth-2\fboxsep-2\fboxrule\relax}
    \begin{tabularx}{\linewidth}{@{}>{\raggedright\arraybackslash\bfseries}p{2.1em} @{\hspace{0.9em}} >{\raggedright\arraybackslash}X @{\hspace{0.9em}} >{\raggedleft\arraybackslash\bfseries}p{2.4em}@{}}
    #1 & #2 & #3
    \end{tabularx}
    \end{minipage}%
  }%
  \par\addvspace{0.32em}%
}


\title{\textbf{Liquid jet adhesion}}
\author{\normalsize X22 (2022) --- English translation}
\date{}
\begin{document}
\maketitle

\section*{Introduction}

\textbf{Attention!} No error calculations are required in this work, except for Part A.

In everyday life, we often observe that a flow of liquid “sticks” to a solid surface. This effect is commonly called the "teapot effect" and occurs in the following processes: inkjet printing and extrusion - producing products by forcing a viscous melt of material or thick paste through a forming hole. In this work, we will quantitatively study the “teapot effect” manifested in the interaction of a cylindrical glass tube and a flow of water, and, as a result, we will obtain a spiral of liquid wrapped around the tube.

\section*{Equipment}

Mariotte vessel with a hose outlet with a nozzle and a tap Tap on a hose Three stands Four legs Four couplings Large container Napkins 6 test tubes of different diameters Sheets of graph paper Sheet of cardboard Protractor Measuring cup Scissors Stopwatch U-shaped pointer Screen with parallel lines Clips

External diameter Internal diameter 1 2.05 1.82 2 1.59 1.37 3 1.37 1.16 4 1.20 1.03 5 0.98 0.81 6 0.96 0.77

\begin{figure}[H]
    \centering
    \includegraphics[width=0.72\textwidth]{E1_EN_assets/fig1.jpg}
    \caption{Figure 1.}
\end{figure}

\section*{Part A. Jet kinematics}

The jet leaving the nozzle is deviated from the straight direction under the influence of gravity. We can use this fact to find the initial flow velocity $U_{0}$. Let us denote the trajectory of water drops as $z(x)$, where $z$ is the vertical coordinate, $x$ is the horizontal coordinate in the plane of the curve.

\problembox{A1}{Experimentally study the $z(x)$ dependence. Take at least 15 measurements. Describe the methodology in detail.}{1.30}

\begin{figure}[H]
    \centering
    \includegraphics[width=0.72\textwidth]{E1_EN_assets/fig2.jpg}
    \caption{Figure 2.}
\end{figure}

\problembox{A2}{Define $U_{0}$.}{0.80}

\problembox{A3}{Determine the internal cross-sectional area $A_0$ of the nozzle and the diameter of the jet at the exit from it $D_j$.}{0.50}

\section*{Part B. Spiral}

Fix the test tube in a vertical position and direct the stream of water as in the figure. Ensure that the jet makes as many revolutions as possible along the surface of the test tube; to do this, you can change the vertical position of the needle to increase or decrease the flow speed and move the test tube so that the flow of water is directed as close to the tangent as possible.

\textbf{Attention!} 1. In parts B and C, the stream should hit the test tube \textbf{horizontally}. 2. In parts B and C, it is necessary to maintain the same initial conditions with high accuracy. As a possible solution to this problem, it is proposed to install the needle horizontally, and also bring it as close as possible to the test tube, which will guarantee a constant value of $U_0$. 3. After measuring the last point, place the U-shaped pointer as shown in the photo below and close the tap. Do not move any of the stands until you have completed all measurements in parts B and C. 4. Laminar flow of water is very important in this experiment, so make sure that at the height you are testing, the stream is laminar for 20 centimeters. This is achieved by selecting the height at which the needle is fixed. 5. The Mariotte vessel is designed to be tight, so as the water flows out, characteristic sounds (gurgling) should be made from time to time. If this does not happen 10 minutes after opening both taps, contact the auditorium attendant.

\begin{figure}[H]
    \centering
    \includegraphics[width=0.72\textwidth]{E1_EN_assets/fig3.jpg}
    \caption{Figure 3.}
\end{figure}

Now temporarily close the tap. For a quantitative description, we introduce the following notation. Let a jet of liquid of density $\rho$ with a volumetric flow rate $Q$ move with a speed $U$ in a spiral wound onto the surface of a vertical cylinder. Let's call the angle between the vertical and the direction of the jet $\psi$. Let the width of the jet be $W$ and its cross-sectional area $A$. Moreover, due to the presence of viscosity inside the jet along its thickness, there is a certain velocity distribution, therefore the speed $U$, which we introduced earlier, is in some way a fictitious value and simply connects the volumetric flow rate $Q$ and the jet area $A$ using the formula $Q = AU$. We describe the spiral using coordinates $(X(s), z(s))$, where $s$ is a natural parametrization, i.e. $ds^2 = dX^2 +dz^2$. In the Cartesian coordinate system, the spiral can be written as
$$\vec{r}(s) = \begin{pmatrix}
R\sin\!\bigl(\frac{X(s)}{R}\bigr) \\
R\cos\!\bigl(\frac{X(s)}{R}\bigr) \\
z(s)
\end{pmatrix}.$$

\begin{figure}[H]
    \centering
    \includegraphics[width=0.72\textwidth]{E1_EN_assets/fig4.jpg}
    \caption{Figure 4.}
\end{figure}

You are invited to experimentally study the dependence $z(X)$. The following method is suggested: make additional marks on a sheet of graph paper to make it easier for you to navigate, then roll it up and insert it into the test tube so that the sheet fits tightly to it.

\begin{figure}[H]
    \centering
    \includegraphics[width=0.72\textwidth]{E1_EN_assets/fig5.jpg}
    \caption{Figure 5.}
\end{figure}

\problembox{B1}{Experimentally obtain the dependence $z(X)$, make at least 30 measurements. Read the warning below.}{1.50}

\textbf{Attention!} 1.) In parts B and C, the stream should hit the test tube \textbf{horizontally}. 2.) Parts B and C need to keep the same initial conditions with high precision. As a possible solution to this problem, it is proposed to install the needle horizontally, and also bring it as close as possible to the test tube, which will guarantee a constant value of $U_0$. 3.) After measuring the last point, place the U-shaped pointer as shown in the photo below and close the tap. Do not move any of the stands until you have completed all measurements in parts B and C. 4.) Laminar water flow is very important in this experiment, so make sure that at the height you are testing, the stream is laminar for 20 centimeters. 5.) The Mariotte vessel is hermetically sealed, so when the taps are open, characteristic sounds should be emitted from the vessel from time to time. If this does not happen 10 minutes after opening both taps, then you need to contact the auditorium duty officer.

\begin{figure}[H]
    \centering
    \includegraphics[width=0.72\textwidth]{E1_EN_assets/fig6.jpg}
    \caption{Figure 6.}
\end{figure}

\problembox{B2}{Plot the dependence $z(X)$.}{0.80}

\problembox{B3}{Calculate $\psi (s)$ and plot this dependence.}{1.30}

\section*{Part C. Dependence of jet shape on test tube radius}

The theory that we will develop in part $D$ predicts a counterintuitive dependence of $z(X)$ on the radius of the test tube $R$ under fixed initial conditions $U_{0}, \psi_0$. However, for the convenience of conducting the experiment, we propose to measure before starting theoretical calculations how $R$ affects the jet parameter $\lambda$ - the height to which water flowing in a spiral descends in one whole revolution. When you try to get a spiral on a test tube of a smaller radius, you may simply get a jet that deviates at some angle. In this case, you can lean the stream closer to the test tube, i.e. ignore the fact that the jet must be as close as possible to touching the test tube.

\begin{figure}[H]
    \centering
    \includegraphics[width=0.72\textwidth]{E1_EN_assets/fig7.jpg}
    \caption{Figure 7.}
\end{figure}

\problembox{C1}{For all test tubes, measure $\lambda (R)$ under the same initial conditions.}{0.80}

\problembox{C2}{Place the points corresponding to the experimental data obtained in the previous paragraph on the graph $z(X)$. Based on this, draw a conclusion regarding the validity of the statement: “The radius of the test tube $R$ does not affect the dependence $z(X)$ under the same initial conditions $(U_{0}, \psi_0)$.”}{0.50}

\section*{Part D. Theoretical calculation of the jet shape}

In this part of the work, all numerical calculations are proposed to be performed using the necessary data from part B.

For further work with the theory, it is proposed to approximate the dependence $\psi (s)$ by a parabola. We will now describe how to do this in the case where the extremum is far from $x = 0$. 1. Let us describe the set of pairs $\{x_i; y_i\}$ by the formula $y = ax^2 + bx + c$. 2. Let us write the deviation function:
$$\sigma(a,b,c,\{x_i,y_i\}) = \sum_i (y(x_i) - y_i)^2 = \sum_i (a x_i^2 + b x_i + c - y_i)^2.$$
3. Set the partial derivatives $\partial\sigma/\partial a$, $\partial\sigma/\partial b$, and $\partial\sigma/\partial c$ to zero. The solution of this linear system minimizes $\sigma(a,b,c;\{x_i,y_i\})$, because $\sigma$ is a positive definite quadratic form in $a$, $b$, and $c$ (clearly $\sigma\to\infty$ if any of $|a|$, $|b|$, $|c|$ grows without bound). 4. Solve the linear system for $a$, $b$, and $c$. The optimality of these coefficients is guaranteed by the minimality of the error function.

\problembox{D1}{Obtain the coefficients $a$, $b$ and $c$ for the parabolic approximation of the $\psi (s) = as^2 + bs + c$ dependence.}{2.00}

To build the theory, we will use the following theses: 1. “Pressing” the water flow to the test tube provides a pressure difference $\Delta P$ between the inner and outer parts of the stream. 2. The jet is acted upon by a viscous friction force with a viscosity coefficient $\eta$. We define viscosity through the force acting on a unit area of the liquid due to the presence of a certain gradient of the parallel component of the velocity of this surface:
$$F_{\text{visc}} = \eta S\,\frac{dv}{dz}.$$
Let's abstract from our spiral to solve a small auxiliary question.

\problembox{D2}{Let us assume that a layer of water of height $H$ flows stationary and uniformly over a flat surface of a solid body. We can assume that the pressure difference between points lying at different ends of the same streamline does not depend on this streamline, because otherwise they would be bent. Find the velocity distribution $u(h)$ inside this layer, where $h$ is the height above the surface. Express your answer in terms of the speed of water at the surface $U_{\text{surf}}$, $h$ and $H$. Express the force of viscous friction $F_{\text{visc},1}$ acting on a layer section of width $W$ and length $L$ from the surface in terms of $U_{\text{surf}}$, $A$, $W$, $L$, $\eta$, where $A = WH$ is the cross-sectional area.}{0.70}

\begin{figure}[H]
    \centering
    \includegraphics[width=0.72\textwidth]{E1_EN_assets/fig8.jpg}
    \caption{Figure 8.}
\end{figure}

We will assume that the distribution of $u(h)$ inside our jet coincides with that described in point B1. We simply take into account the difference in the geometry of the section from a rectangle in the form of a constant dimensionless factor $C$, i.e. if you got the answer $F_{\text{visc},1} = f(U_{\text{surf}}, H, W, L, \eta)$ in D2, then a force $F_{\text{visc},2} = C \cdot f(U_{\text{surf}}, H, W, L, \eta)$ acts on a section of the jet of length $L$. We will also assume that the width of the jet $W$ does not change when moving along the spiral and is equal to the diameter of the jet $W = D_j$ at the exit from the source. \textit{Note:} If you did not do point A3, then take $D_j = (2.00 \pm 0.04)~\text{mm}$ as known.

\problembox{D3}{Express $U_{\text{surf}}$ in terms of $U$. \textit{Recall} that $U$ is some characteristic quantity determined through the expression $Q = AU$.}{0.60}

\begin{figure}[H]
    \centering
    \includegraphics[width=0.72\textwidth]{E1_EN_assets/fig9.jpg}
    \caption{Figure 9.}
\end{figure}

In paragraphs D4 - D8 we will describe the dynamics of a section of a water spiral of length $ds$.

\problembox{D4}{Formulate the law of conservation of mass for this piece. Express your answer in terms of $U$, $A$, $ds$, $\rho$.}{0.20}

\problembox{D5}{Write Newton's vector II law in impulse form. Take as a basis the orthonormal triple of vectors $\left(\vec{t} (s), \vec{n} (s), \vec{b} (s) \right)$: $\vec{t}$ is directed along the jet; $\vec{n}$ - directed perpendicular to the surface of the test tube from it; $\vec{b}$ is their orthogonal complement (in the tangent plane perpendicular to $\vec{t}$). When recording the change in the momentum of the jet, neglect the distribution of velocities inside it and consider it a uniform flow over the cross section, moving at a speed of $U(s)$, this approximation can be used due to the fact that taking into account the unevenness of the velocity in this term gives a correction very close to 1. Use point $D4$ to simplify this expression.}{1.50}

\problembox{D6}{Eliminate time from the equation using the identity of the operators
$$\frac{d}{dt} \equiv U\frac{d}{ds}.$$}{0.20}

\problembox{D7}{Find the projections of $\cfrac{d}{ds} ( U \vec{t} ) $ onto $\vec t$ and $\vec{b}$. Express your answers in terms of functions of $U$, $\psi$ and their derivatives with respect to $s$.}{0.60}

\problembox{D8}{Using the results of points D4--D7, obtain a system of two differential equations for $U$ and $\psi$ of the following form:
$$\begin{cases}
\cfrac{dU}{ds} = -\alpha_1 U^2 + \cfrac{\alpha_2\cos\psi}{U}, \\[6pt]
\cfrac{d\psi}{ds} = -\cfrac{\beta\sin\psi}{U^2},
\end{cases}$$
where $\alpha_1$, $\alpha_2$, and $\beta$ are positive constants. Express them in terms of $\eta$, $U_0$, $D_j$, $\rho$, $g$, and $C$.}{0.60}

\problembox{D9}{Calculate the value of $U$ at the experimental points based on the parabolic approximation $\psi (s)$ (point D1). Error $U$ may not be assessed. Graph $U(s)$.}{1.20}

\problembox{D10}{Rate the value $C$.}{0.80}

To check the obtained value $C$, we numerically solve the system of equations obtained in step D8 iteratively:
$$\begin{cases}
U(s+\Delta s) \approx U(s) + \cfrac{dU}{ds}\,\Delta s, \\[6pt]
\psi(s+\Delta s) \approx \psi(s) + \cfrac{d\psi}{ds}\,\Delta s.
\end{cases}$$

\problembox{D11}{Solve differential equations approximately numerically. Find the type of “theoretical” dependence $U_t(s)$, $\psi_t(s)$, $X_t(z)$ at at least 15 points.}{1.20}

\problembox{D12}{On the graphs $z(X)$, $\psi (s)$, $U(s)$, mark the theoretical modeling points $z_t(X)$, $\psi_t(s)$, $U_t(s)$.}{0.90}

\section*{Part E. Spiral Observability}

We will study the observability of the spiral depending on the initial flow velocity $U_0$. Let's consider a situation when the flow speed is so high that a spiral does not form, and the water flow simply deviates by some angle $\alpha$ from the initial direction. Try to ensure that the flow of water when touching the test tube is as close to horizontal as possible, and its initial direction is as close to the tangent as possible.

\begin{figure}[H]
    \centering
    \includegraphics[width=0.72\textwidth]{E1_EN_assets/fig10.jpg}
    \caption{Figure 10.}
\end{figure}

\problembox{E1}{Experimentally study the dependence $\alpha (U_0)$ and construct its graph. Take at least 13 measurements. Describe the methodology in detail.}{1.50}

\problembox{E2}{At what maximum $U_{\mathrm{crit}}$ is a spiral observed?}{0.50}

\section*{Constants}

$\rho = 1.00~\text{kg}/\text{m}^3$, $\eta = 8.9 \cdot 10^{-4}~\text{Pa}\cdot\text{s}$, $g = 9.8~\text{m}/\text{s}^2$

\vspace{1em}
\noindent\textit{Source:} \url{https://pho.rs/p/507}

\end{document}